\chapter{Fonts}

OmniLaTeX provides a comprehensive font system built on LuaLaTeX's native OpenType
support. All fonts are configured through a layered architecture that separates
font selection, feature configuration, and fallback logic.

% =========================================================================
\section{Font System Architecture}
% =========================================================================

OmniLaTeX targets LuaHBTeX 1.21.0+, which provides native OpenType and
Unicode support without the shims required by XeLaTeX. The font stack is
composed of three core packages:

\begin{description}
  \item[\texttt{fontspec}] Provides \cs{setmainfont}, \cs{setsansfont},
    \cs{setmonofont}, and raw OpenType feature control.
  \item[\texttt{unicode-math}] Replaces the legacy \texttt{amsmath} math font
    interface with OpenType math fonts.  Loads \cs{symup}, \cs{symbf},
    \cs{symit}, and the full \textsc{symb} alphabet.
  \item[\texttt{lualatex-math}] Fixes a handful of LuaTeX math-mode bugs
    (e.g.\ \cs{sqrt} spacing, \cs{overline} thickness) that affect
    LibreOffice and other engines.
\end{description}

Font definitions live in the \texttt{omnilatex-fonts} module and are loaded
before the user preamble, so every document inherits the defaults
automatically.

% =========================================================================
\section{Default Font Stack}
% =========================================================================

\begin{table}[htbp]
  \centering
  \caption{Default font families shipped with OmniLaTeX.}
  \label{tab:default-fonts}
  \begin{tabular}{@{}lll@{}}
    \toprule
    Family        & Default Font               & Module            \\
    \midrule
    Roman (serif) & Libertinus Serif           & \texttt{omnilatex-fonts} \\
    Sans-serif    & Libertinus Sans            & \texttt{omnilatex-fonts} \\
    Math          & Libertinus Math            & \texttt{omnilatex-fonts} \\
    Monospaced    & Monaspace Neon             & \texttt{omnilatex-fonts} \\
    Accessibility & Atkinson Hyperlegible Next & \texttt{omnilatex-fonts} \\
    \bottomrule
  \end{tabular}
\end{table}

Libertinus is the spiritual successor to Linux Libertine and provides a
complete serif, sans, and math family with consistent metrics.
Monaspace Neon offers five weight-variant glyphs in a single monospaced
face, making it ideal for code listings.  Atkinson Hyperlegible Next is
loaded as an accessibility fallback for readers with low vision.

% =========================================================================
\section{Setting Fonts}
% =========================================================================

OmniLaTeX exposes four high-level commands that wrap the underlying
\texttt{fontspec} API:

\begin{table}[htbp]
  \centering
  \caption{Font-setting commands.}
  \label{tab:font-commands}
  \begin{tabular}{@{}ll@{}}
    \toprule
    Command                          & Underlying call            \\
    \midrule
    \texttt{\textbackslash setMainFont\{...\}}     & \texttt{\textbackslash setmainfont}  \\
    \texttt{\textbackslash setSansFont\{...\}}     & \texttt{\textbackslash setsansfont}  \\
    \texttt{\textbackslash setMonoFont[...]\{...\}} & \texttt{\textbackslash setmonofont}  \\
    \texttt{\textbackslash setMathFont\{...\}}     & \texttt{\textbackslash setmathfont}  \\
    \bottomrule
  \end{tabular}
\end{table}

Minimal example that switches to Source Serif~4 and Source Code~Pro:

\begin{verbatim}
\usepackage{omnilatex-fonts}
\setMainFont{Source Serif 4}
\setSansFont{Source Sans 3}
\setMonoFont[Scale=0.88]{Source Code Pro}
\setMathFont{STIX Two Math}
\end{verbatim}

% =========================================================================
\section{Font Feature Tables}
% =========================================================================

OpenType fonts expose named features (e.g.\ \texttt{liga}, \texttt{kern},
\texttt{onum}, \texttt{ss01}) that control ligatures, kerning, glyph
substitution, and stylistic alternates.  OmniLaTeX enables sensible
defaults for every family:

\begin{table}[htbp]
  \centering
  \caption{Default OpenType features by font family.}
  \label{tab:ot-features}
  \begin{tabular}{@{}lcccc@{}}
    \toprule
    Feature & Roman & Sans & Mono & Math \\
    \midrule
    \texttt{kern}   & on & on & on & --- \\
    \texttt{liga}   & on & on & on & --- \\
    \texttt{onum}   & on & on & on & --- \\
    \texttt{pnum}   & on & on & on & on  \\
    \texttt{ss01}   & on & off & off & --- \\
    \bottomrule
  \end{tabular}
\end{table}

To query available features for any installed font, run:

\begin{verbatim}
lualatex --luaonly -e "
  local f = require'fontloader.font'
  -- inspect font's 'features' table
"
\end{verbatim}

Or use the \texttt{otfinfo} command-line tool bundled with \TeX\ Live:

\begin{verbatim}
otfinfo --features LibertinusSerif-Regular.otf
\end{verbatim}

To override features per family, pass the \texttt{Features} key:

\begin{verbatim}
\setMainFont{Libertinus Serif}[
  RawFeature={+ss01,+salt=3},
  Ligatures=TeX,
]
\end{verbatim}

% =========================================================================
\section{Icon Fonts}
% =========================================================================

OmniLaTeX loads \texttt{fontawesome5}, which provides over 1\,500 scalable
icons as font glyphs.  The primary interface is \cs{faIcon}\texttt{\{name\}}:

\begin{verbatim}
\faIcon{graduation-cap}  \faIcon{github}  \faIcon{python}
\end{verbatim}

\begin{table}[htbp]
  \centering
  \caption{Commonly used Font Awesome icons.}
  \label{tab:fa-icons}
  \begin{tabular}{@{}lll@{}}
    \toprule
    Command                          & Icon & Use case          \\
    \midrule
    \texttt{\textbackslash faIcon\{graduation-cap\}} & \faIcon{graduation-cap} & Education \\
    \texttt{\textbackslash faIcon\{github\}}          & \faIcon{github}          & Version control \\
    \texttt{\textbackslash faIcon\{python\}}          & \faIcon{python}          & Programming \\
    \texttt{\textbackslash faIcon\{envelope\}}        & \faIcon{envelope}        & Contact \\
    \texttt{\textbackslash faIcon\{home\}}            & \faIcon{home}            & Navigation \\
    \texttt{\textbackslash faIcon\{book\}}            & \faIcon{book}            & References \\
    \texttt{\textbackslash faIcon\{warning\}}         & \faIcon{warning}         & Alerts \\
    \texttt{\textbackslash faIcon\{lightbulb\}}       & \faIcon{lightbulb}       & Tips \\
    \bottomrule
  \end{tabular}
\end{table}

Icons inherit the surrounding text colour and scale, so they integrate
seamlessly into headings, footers, and inline text.

% =========================================================================
\section{CJK Fonts}
% =========================================================================

Chinese, Japanese, and Korean (CJK) typesetting is handled by
\texttt{luatexja-fontspec}, a bridge between \texttt{luatexja} and
\texttt{fontspec}.  OmniLaTeX provides three convenience commands:

\begin{verbatim}
\setCJKMainFont{Noto Serif CJK SC}
\setCJKSansFont{Noto Sans CJK SC}
\setCJKMonoFont{Noto Sans Mono CJK SC}
\end{verbatim}

The \texttt{Noto CJK} superfamily covers four regional variants:

\begin{table}[htbp]
  \centering
  \caption{Noto CJK regional variants.}
  \label{tab:cjk-variants}
  \begin{tabular}{@{}lll@{}}
    \toprule
    Tag & Script / Region     & Example font              \\
    \midrule
    SC  & Simplified Chinese  & Noto Serif CJK SC         \\
    TC  & Traditional Chinese & Noto Serif CJK TC         \\
    JP  & Japanese            & Noto Serif CJK JP         \\
    KR  & Korean              & Noto Serif CJK KR         \\
    \bottomrule
  \end{tabular}
\end{table}

See Chapter~\ref{ch:cjk} for full CJK typesetting details including
punctuation handling, ruby annotations, and inter-paragraph spacing.

% =========================================================================
\section{RTL Fonts}
% =========================================================================

Right-to-left (RTL) scripts such as Arabic, Hebrew, and Persian require
bidirectional text shaping.  OmniLaTeX uses \texttt{bidi} (LuaLaTeX) with
font fallback chains:

\begin{verbatim}
\setArabicFont{Noto Naskh Arabic}[Script=Arabic]
\setHebrewFont{Noto Serif Hebrew}[Script=Hebrew]
\end{verbatim}

The \texttt{bidi} package automatically mirrors punctuation and adjusts
mixed LTR/RTL paragraph direction.  Fallback chains ensure that glyphs
missing from the primary font are drawn from a secondary source.

See Chapter~\ref{ch:rtl} for advanced RTL configuration including
mixed-script documents and table direction.

% =========================================================================
\section{Unit Number Font}
% =========================================================================

Scientific documents that use \texttt{siunitx} benefit from a dedicated
number font that guarantees consistent digit width and alignment.
OmniLaTeX defines \cs{unitnumberfont} for this purpose:

\begin{verbatim}
\sisetup{
  number-font = \unitnumberfont,
}
\end{verbatim}

This selects the monospaced face (Monaspace Neon by default) so that
tabular number columns align on the decimal point without requiring
\texttt{S} column type tricks.  The command can be redefined to use any
font with tabular figures:

\begin{verbatim}
\renewcommand{\unitnumberfont}{\fontspec{Fira Code}}
\end{verbatim}
